\section{Constructions}\label{sec:constructions}

\paragraph{The unit.} A construction pairs a \emph{form} with a \emph{meaning} and an \emph{effect}. The form is a
pattern over what is said. The meaning is what a reader obtains when the form matches. The effect is an affordance:
what the construction makes available when its meaning is acted on, such as a record that may be written to memory
or a step that an agent is licensed to take. A construction that matches does not act; it makes an action
available. The pairing of form and meaning is the central idea of construction grammar, which Fluid Construction
Grammar (FCG) makes operational \citep{beuls2023fcgstate}. The effect is the addition of the companion paper
\citep{vera2026constructions}.

Constructions exist at every level of generality \citep{beuls2023cxgai}: a coined term, a fixed phrase that stands
for a settled argument, a pattern with open slots such as ``$X$ is a coordination problem, not a technical one'',
or an abstract schema for stating a claim with its evidence. A finite inventory of such units combines to read and
produce an unbounded range of expressions.

One construction that exists today shows the three parts \citep{vera2026constructions}. A member teaches the phrase
``file $X$'' to mean ``make a ticket from $X$''. The definition is compiled, with no model, into a construction
whose form is the word \emph{file} followed by a slot, whose meaning is a one-step procedure with the slot's text
as its title, and whose effect is that the agent may run that procedure. When someone later says ``file broken
printer'', the reading names the procedure and its title, and removing the construction leaves the sentence without
that reading.

Three properties make the construction a workable unit for a protocol. It can be \emph{individuated}: it can be
named, hashed and counted. It can be \emph{attributed}: its first registration has a time and an origin. And its use
is an \emph{event} that anyone holding the construction and the text can repeat.

\paragraph{A symbolic harness.} An agent built on a language model is governed by a harness of prompts, tool
definitions and control code that engineers write. A
set of constructions does the same work. It reads each turn, and what it reads determines which steps are available
to the agent. The engine's own search decides which construction applies, with no model in the loop, and the
protocol requires that the same turn yield the same affordances every time (Section~\ref{sec:encounters}). A harness written in constructions can be taught in conversation,
read, attributed, and served to anyone's agents. That is why it can be published.

\paragraph{The grammar as a product.} Constructions are of two kinds. The \emph{base grammar} reads ordinary
language and is the same everywhere; it is nobody's product and earns nothing. A Bonfire's \emph{own}
constructions are those its members taught or derived from its discourse. A \emph{grammar}, in this paper, is a
versioned set of a Bonfire's own constructions that it chooses to publish. The grammar is the Bonfire's product: a
working capability, published for anyone to use. With it the Bonfire publishes a \emph{signal}, a statement of what
the grammar is for, by which others find it (Section~\ref{sec:adoption}).

The grammar is the product, but the construction is the unit of account. Credit (Section~\ref{sec:reward}) is
counted per construction and never per grammar. If it were counted per grammar, a publisher could fill a grammar
with constructions that do nothing and be credited for the bundle.

\paragraph{Standing.} A Bonfire's \emph{inventory} is everything its agents read with: its own grammar, the constructions it has
taken from others (Section~\ref{sec:adoption}), and the base grammar. Inside
an inventory each construction has a \emph{standing}, a score that records how well it has served there. We adopt
the rule published for language-game experiments in FCG \citep{nevens2022language}: a new construction starts at
$0.5$; on a successful use the constructions used gain $0.1$ and their competitors lose $0.1$; on a failed use the
constructions used lose $0.4$; and a construction whose score reaches $0$ is deleted. The authors report that,
within conditions they state, the exact values do not determine the overall dynamics of learning. The shape is what
matters here: success reinforces, competitors are inhibited, and failure costs more than success earns. Standing
is local and is never paid on. Its one job in the protocol is to decide what an inventory keeps
(Section~\ref{sec:adoption}).
